% Draft. acmart/sigconf, verified to build on this box on 2026-08-13
% after collection-latexextra completed; see paper/README.md for how the
% environment was checked (by compiling, not by kpsewhich).
%
% To fall back to a class that needs nothing installed, swap the two
% lines below and drop \acmConference..\keywords. The article version
% built 3 pages with no warnings and is what every claim here was
% drafted against.
\documentclass[sigconf,nonacm]{acmart}
% \documentclass[10pt,twocolumn]{article}\usepackage[margin=0.75in]{geometry}

\usepackage{booktabs}
\usepackage{microtype}
\usepackage{amsmath}

\newcommand{\ci}[2]{[\,#1,\ #2\,]}

\title{Your peer sets your step time: spatial GPU partitioning is
pair-specific, not width-specific}
\author{Draft}
\affiliation{\institution{Claims traced to
  \texttt{docs/claims-and-evidence.md}}\country{}}

\begin{document}

\begin{abstract}
We partition a single AMD GPU between concurrent diffusion tenants using
CU masks and a step-granular scheduler. Any such scheduler needs to know
what a step costs at $q$ units and how much more it costs when a peer
holds the rest. We report that the second quantity is not a function of
the split at all.

\textbf{It is set by which peer.} On gfx1201, SDXL beside CogVideoX-2b
pays $1.689\times$ its solo step at $4{+}28$ and $26.137\times$ at
$28{+}4$ --- the penalty rises as the image tenant's own slice
\emph{widens}, because its co-run step obeys $\mathit{corun}_{\text{sdxl}}
= \mathit{corun}_{\text{cog}} + 0.54\,\mathit{solo}_{\text{sdxl}}$, with
$k \in [0.531, 0.560]$ over six runs. The dominant term is the video
tenant's step, which the image tenant's width does not change. Fitted on
three splits and tested out of sample on the two not yet run, predicting
$1736$ and $3369$\,ms against $1659.7$ and $3213.2$ measured. Going from
4 units to 28 improves the image tenant's solo step $4.1\times$ and its
co-run step $3.8\times$ \emph{worse}. It is one-directional: the video
tenant pays $1.006$--$1.064$ beside SDXL and $1.276$ beside another
CogVideoX. Against rotation --- the whole die to each tenant half the
time --- the pair's aggregate step rate is $5.608$/s against
$0.629$--$2.889$/s partitioned, so the best split is $1.94\times$ worse
and the worst $8.9\times$.

\textbf{It counts peers, not busy die.} A slice pays $1.050$, $1.367$,
$2.318$, $4.520$ at 1, 2, 4 and 8 equal ways on gfx1201 and $1.024$,
$1.247$, $1.507$, $2.245$ on gfx90a. Four ways is $+77.3\%$ over the
pairwise entry for a slice of the same width.

\textbf{It stops at the process boundary.} Two processes on disjoint
$16{+}16$ masks pay $1.97\times$ each; giving both the whole die instead
costs $2.10\times$. Two \emph{threads} of one process on those same
halves pay $1.37\times$. In aggregate step rate, masking across processes
costs $24\%$ against not masking, so the mechanism is available only to a
runtime holding every tenant in one address space.

\textbf{And for same-model tenants it is drawn, not computed.} The
penalty at $16{+}16$ is bistable: either $1.297\times$ or $1.949\times$,
latched per process, identifiable from a single step, slow in $16.9\%$,
\ci{11.4\%}{23.6\%} of 160 processes. It is absent on gfx90a.

The scheduling consequence is negative and pre-registered. On the
mismatched workload our grid was designed around --- 405 cells, nine
policies --- the method is $0.56\%$ \emph{worse} than the strongest
baseline against a $20\%$ target. Strict priority, which never partitions
at all, beats the adaptive scheduler on urgent miss rate by $+134\%$
under arrivals and $+84.5\%$ under backlog. A dual-ledger probe that
appeared to cut urgent miss by $13\%$ was avoiding a conservative
envelope we had added ourselves; with the envelope removed its benefit is
$+0.41\%$, \ci{-0.0211}{+0.0256}.

Finally, a methodological result we would report even if the rest had
gone the other way. Three separate cost models here were beautiful,
quantitative and wrong, each an artifact of a deferred-event instrument
that repeats one stale reading when the CPU runs ahead --- including one
fitting five splits to within $3$--$9\%$ over a $3\times$ range. The
control that catches all three is cheap and was missing from every one:
re-measure the solo \emph{after} the co-run, with the peers idle. A
penalty that does not vanish when the peers stop is not a penalty.
\end{abstract}

\maketitle

\section{Introduction}

Serving a latency-sensitive image request beside a long-running video
generation on one GPU is a scheduling problem with an unusual property:
the two tenants can be given \emph{disjoint sets of compute units}
rather than alternating slices of time. On AMD hardware this is
\texttt{hipExtStreamCreateWithCUMask}; on NVIDIA it is the SM mask.
Spatial partitioning is attractive because neither tenant has to be
suspended, and a suspended diffusion step is not free to resume.

The premise of any such scheduler is a cost model: how much does a step
cost at $q$ units, and how much more does it cost when a peer holds the
other $32-q$? This paper is about the second question, and our answer is
that it is not a question about $q$. The externality is a property of the
\emph{pair} --- which models, how many of them, and whether they share an
address space --- and for a mismatched pair it is dominated by a term the
tenant's own width cannot move. A scheduler that models it as
$f(q, 32-q)$, as ours did and as a cost-model-driven partitioner must,
is not approximating this hardware; it is describing a different one.

\paragraph{Contributions.}
\begin{enumerate}
\item A measurement of the co-run penalty on gfx1201 showing it is
  bistable and latched per process --- drawn rather than computed from
  the split (\S\ref{sec:bistable}).
\item The paper's central result: for a mismatched pair the co-run step
  time is set by the \emph{slower} tenant, so widening the faster
  tenant's slice makes it worse, and every split loses to rotation. Five
  splits on gfx1201, pre-registered, with the law tested out of sample
  (\S\ref{sec:mismatch}).
\item Two further ways the penalty escapes $f(q, 32-q)$: it counts peers
  rather than busy die, at 1, 2, 4 and 8 ways on two architectures; and
  it does not cross the process boundary, so the mechanism exists only
  inside one address space (\S\ref{sec:mismatch}).
\item Three negatives. Two are about our own design: a probe whose
  measured benefit is entirely the avoidance of a conservative envelope
  we added (\S\ref{sec:probe}), and a pre-registered failure on the
  workload the grid was built around (\S\ref{sec:negative}). The third is
  that strict priority --- no partitioning at all --- beats the adaptive
  scheduler on this workload.
\item A methodological account of three cost models that were beautiful,
  quantitative and wrong for the same instrument reason, and of the one
  control that catches all three (\S\ref{sec:method}).
\end{enumerate}

\paragraph{What this paper does not claim.} That spatial partitioning is
useless. Masking is exact here, byte-for-byte reproducible across
re-grants, and beats two unmasked full-die streams. What we claim is that
its cost cannot be predicted from widths, that on a mismatched pair it
loses to time-slicing, and that a scheduler must therefore keep temporal
isolation as a first-class action and refuse configurations it has no
measurement for.

\section{Setup}

One Radeon AI PRO R9700 (gfx1201, 32 maskable units, 34.2\,GB). SDXL at
$768\times768$; CogVideoX-2b at $480\times720$, 9 frames. A step-granular
executor exposes \texttt{prepare}, \texttt{run\_step}, \texttt{suspend},
\texttt{resume} and \texttt{finalize}; suspension captures the latent,
the step index, the scheduler's internal cursor and the generator state.
Interrupting a request between every step and re-granting it at a
different mask width produces a latent identical, byte for byte, to the
uninterrupted baseline's.

Masks are installed per stream and read back every time. A stream that
quietly carried the whole die would produce an unusually \emph{low}
externality, which reads as good news, so the readback is evidence
rather than a sanity check.

\paragraph{Measuring a step.}
Step cost is device time between two CUDA events bracketing the step on
its own stream. Reading those events with a per-step
\texttt{synchronize} drains the pipeline and charges the drain to the
measurement: it made a $0.1155$\,s step read as $0.256$\,s. Reading them
one step late avoids the drain but introduces a failure discussed in
\S\ref{sec:method}.

\section{The co-run penalty is bistable}
\label{sec:bistable}

Two SDXL tenants on disjoint $16{+}16$ masks pay one of two costs.

\begin{table}[t]
\centering
\begin{tabular}{lrrr}
\toprule
state & processes & mean & range \\
\midrule
fast & 23 of 30 & $1.297$ & $1.287$--$1.303$ \\
slow & 7 of 30 & $1.949$ & $1.923$--$1.962$ \\
\bottomrule
\end{tabular}
\caption{Co-run cost as a multiple of the solo step, one campaign.
The clusters do not overlap.}
\end{table}

Three properties make this actionable rather than merely annoying.

\textbf{It latches.} No process was observed in both states: 0 of 8
processes changed across nine later episodes each, although every
episode built fresh adapters and re-acquired its streams. Re-forming a
pairing does not redraw.

\textbf{It is per mask pair.} Interposing an $8{+}24$ pairing never
disturbed the $16{+}16$ state (16 of 16 checks), and $8{+}24$ carried its
own independent draw.

\textbf{One step identifies it.} Comparing each episode's first step
against its own median gives a mean absolute error of $0.00\%$ over 72
episodes; the per-step series inside an episode is constant.

Pooled over 160 processes the slow state occurs in 27, $16.9\%$,
Clopper--Pearson \ci{11.4\%}{23.6\%}. The rate is \emph{not} stable
across campaigns --- $22.7\%$, $35.0\%$, $8.3\%$, $11.7\%$, with Fisher
$p=0.025$ between the extremes --- and we report no predictor. Five were
proposed and retracted: working set, power cap, uncontrollable
bistability, launch stagger, and prior use of a full-die mask.

\paragraph{It is RDNA4's, not AMD's.} On an MI250X (gfx90a, CDNA2, 104
units) the same measurement --- same instrument, same model, the same
half-and-half split at $52{+}52$, 32 processes rotating across four GCDs
--- gives $1.2336$ with a standard deviation of $0.0030$ and nothing
above the slow threshold. At the gfx1201 rate that has probability
$0.0027$. The four GCDs do differ from one another, deterministically, by
$0.6\%$ with non-overlapping ranges, so the instrument resolves far less
than a $50\%$ bimodality on that machine. Partitioning still pays there,
$+5.6\%$ against rotation. The mechanism generalises; the pathology does
not.

The consequence for a cost model is direct. Against rotation at
$2\times108.9$\,ms, the fast state makes partitioning worth $+11.5\%$ and
the slow state worth $-25.5\%$. A model with one externality per width
pair is calibrating against a draw it cannot see.

\section{The slower tenant sets the pace}
\label{sec:mismatch}

\begin{table}[t]
\centering
\begin{tabular}{lrrrrr}
\toprule
split & SDXL & SDXL & SDXL & same & Cog \\
(sdxl+cog) & solo & co-run & ext & model & ext \\
\midrule
$4{+}28$ & $506.2$ & $855.1$ & $1.689$ & $1.139$ & $1.064$ \\
$8{+}24$ & $262.1$ & $780.0$ & $2.975$ & $1.235$ & $1.057$ \\
$16{+}16$ & $153.7$ & $920.1$ & $5.987$ & $1.324$ & $1.050$ \\
$24{+}8$ & $124.0$ & $1659.7$ & $13.386$ & $1.202$ & $1.025$ \\
$28{+}4$ & $122.9$ & $3213.2$ & $26.137$ & $1.115$ & $1.006$ \\
\bottomrule
\end{tabular}
\caption{Milliseconds per denoising step, wall clock with a per-adapter
drain, last of six episodes. \emph{same model} is SDXL beside SDXL at the
identical width, measured in the same session on the same instrument.
Every control returns to within $3.2\%$ of the pre-episode solo and every
solo is within $3\%$ of the measured quota curve.}
\label{tab:mismatch}
\end{table}

Table~\ref{tab:mismatch} is the paper's central measurement. Read down
the SDXL column: the image tenant's slice grows from 4 units to 28, its
solo step improves $4.1\times$, and its co-run step gets $3.8\times$
\emph{worse}. Against its own same-model control at the identical width
it pays $1.48\times$ at $4{+}28$ and $23.43\times$ at $28{+}4$.

The reason is visible in the two co-run columns. Across all six
mismatched runs, including one with the offsets reversed,
%
\begin{equation}
\mathit{corun}_{\text{sdxl}} = \mathit{corun}_{\text{cog}}
  + k\,\mathit{solo}_{\text{sdxl}},\qquad k \in [0.531, 0.560],
\label{eq:pacing}
\end{equation}
%
mean $0.543$, spread $5.5\%$ of it. The dominant term is the video
tenant's step. The image tenant's width does not appear in it, and
widening the image tenant's slice enters only by narrowing the video
tenant's, which makes the dominant term larger: $\mathit{corun}_{\text{cog}}$
runs from $581$\,ms at 28 units to $3148$\,ms at 4. More die, for this
tenant, on this pair, is strictly worse.

Equation~\ref{eq:pacing} was fitted on the first three splits to land and
written down, with its predictions for the two not yet run, before those
ran: $1736$ and $3369$\,ms against $1659.7$ and $3213.2$ measured, $4.4\%$
and $4.6\%$ out. We record this because a five-point fit to five points
is not evidence of a law, and because this project has published a
beautiful five-split fit before that was fitting an instrument
(\S\ref{sec:method}).

\paragraph{It is one-directional, and it follows the model.} CogVideoX
pays $1.006$--$1.064$ beside SDXL and $1.276$ beside another CogVideoX at
the same width, so the mismatch is $0.82\times$ \emph{cheaper} for it.
Nothing here is symmetric contention: the short-step tenant pays the
long-step tenant's step. Reversing the offsets at $16{+}16$, so the
models exchange halves of the die, moves the penalty by $0.06$ --- it
travels with the tenant, not with the position.

\paragraph{Against rotation.} A side beats rotation --- the whole die for
half the wall clock --- iff its co-run step is under twice its whole-die
solo. The image tenant loses at all five splits. The video tenant wins at
$4{+}28$, $8{+}24$ and $16{+}16$ and loses at the two where its own slice
is narrow. For the pair together, aggregate step rate is $5.608$/s under
rotation against $2.889$, $2.859$, $2.279$, $1.230$ and $0.629$/s
partitioned: the best split $1.94\times$ worse, the worst $8.9\times$.

\paragraph{Blocked or slowed, we do not say.} The image tenant's
event-timed step tracks its own solo at every split while its wall clock
reads up to $26\times$ that; the video tenant's two instruments agree to
$0.1\%$. Either the fast tenant's deferred event reading is stale --- the
failure of \S\ref{sec:method}, and the fast tenant is exactly the one
whose CPU races --- or its kernels are unslowed and its thread is blocked
between steps. The throughput result is the same either way, 13 steps in
$12.0$\,s against $2.0$\,s alone, but the mechanism is not, and our
stored series cannot separate them.

\section{The penalty counts peers, and stops at the process boundary}
\label{sec:peers}

Two more measurements say the externality is not $f(q, 32-q)$, and
neither needs a mismatched pair.

\paragraph{Peers, not busy die.} With $N$ equal slices of one model
covering the whole die, a slice pays $1.050$, $1.367$, $2.318$ and
$4.520$ at $N = 1, 2, 4, 8$ on gfx1201, and $1.024$, $1.247$, $1.507$ and
$2.245$ on gfx90a. In every case the die is fully subscribed and only the
number of peers changes. Four slices of 8 units pay $77.3\%$ more than
the pairwise table's entry for an 8-unit slice with one 24-unit peer. The
$N = 2$ cell is the control that licenses the rest: there the
arrangement \emph{is} the pairwise one, and it agrees with the
independently measured pairwise entry to $2.4\%$.

\paragraph{Not across processes.} Every co-run number above, and every
one we know of in this line of work, is measured with threads of one
process. Deployed multi-tenant serving is normally one server per tenant.
Two processes on disjoint $16{+}16$ masks pay $1.97\times$ each; giving
both processes the whole die instead costs $2.10\times$ --- close enough
to $2\times$ to be the hardware scheduler time-slicing them. The mask
still applies, verified against the quota curve at $144.6$ versus
$112.8$\,ms unmasked, so it caps each process at half the units and the
time-slicing happens anyway. In aggregate step rate: one process alone on
the whole die $8.87$/s, two threads of one process at $16{+}16$
$9.54$/s, two processes with the whole die each $8.45$/s, two processes
on disjoint halves $6.46$/s. Masking across processes costs $24\%$ and
buys nothing.

\section{A scheduler that measures, and what it was really doing}
\label{sec:probe}

The runtime charges each tenant from the \emph{measured} cost of the step
it just ran, never from the cost model. Charging from prediction would
make the accounting exact by construction and blind. The policy pairs
tenants by matched step cost and probes: if a formed pairing's observed
step exceeds its paired expectation by a tolerance, the pairing is
dropped for a round. The runtime adds a safety envelope that falls back
to a serial action when a pairing has no measured profile or when drift
exceeds 15\%.

Decomposed with three arms from one code path --- pairing only, plus the
deadline actions, plus the probe --- paired by seed, the probe carried
$-13.14\%$ in the fast state and $-14.20\%$ in the slow one, with the
deadline actions null in both, and it cut serial fallbacks by 73--81\%.

\begin{table}[t]
\centering
\begin{tabular}{lrr}
\toprule
policy & envelope on & envelope off \\
\midrule
static SM partition & $0.3011$ & $0.2415$ \\
EDF & $0.2999$ & $0.2574$ \\
step-matched pairing & $0.3011$ & $0.2415$ \\
probe & $0.2619$ & $0.2353$ \\
\bottomrule
\end{tabular}
\caption{Urgent SLO miss rate, 9 paired seeds. Every policy improves
with the envelope off; all four intervals exclude zero.}
\end{table}

\textbf{Then we removed the envelope.} Every policy improved. And the
probe's advantage over step-matched pairing went from $-13.01\%$,
\ci{-0.0668}{-0.0132}, to $+0.41\%$, \ci{-0.0211}{+0.0256} --- an
interval spanning zero with the point estimate on the wrong side.

So the mechanism was identified correctly and its significance was
backwards. The probe does cut serial fallbacks; those fallbacks are
harmful; and they are harmful because they are ours. The envelope fired
on drift measured as $|{\text{observed}/\text{expected}}-1|$, which
treats a pairing arriving \emph{cheaper} than predicted as a reason to
take the whole die serial. Correcting it to fire on over-run only removed
half the firings, and removing it entirely was better still. The
envelope's other branch, a fallback when the cost model has no entry for
a pairing, never fired in any workload we ran, so it is untested rather
than refuted.

The best configuration we measured is the envelope off and no probe:
step-matched pairing at $0.2276$ against $0.2286$ with the probe.

\paragraph{Robustness.} Injecting predictor error into the belief the
policy reads, and nowhere else, gives zero safety failures at $\pm20\%$
across 134 cells with the envelope both on and off, where a safety
failure is defined in code as over-committing the die, charging from the
cost model when a measurement existed, or charging a measurement taken at
another quota.

\section{A pre-registered negative}
\label{sec:negative}

On the mismatched workload --- 405 cells, nine policies, five seeds,
loads $\{0.6,0.85,1.05\}$ and bursts $\{2,4,8\}$ --- the method is
$0.56\%$ \emph{worse} than the strongest non-oracle baseline,
\ci{+0.0000}{+0.0056} over 45 paired configurations, against a target of
a 20\% improvement.

It follows from \S\ref{sec:mismatch}, though not in the way we first
wrote it down. We originally explained it by saying mismatched pairings
cost $1.02$--$1.06$, so a contention manager has nothing to manage. That
explanation was built on a measurement we have since withdrawn. The
replacement runs the other way: partitioning a mismatched pair is not
cheap but ruinous for the faster tenant, so a policy whose action is to
partition is choosing the worse arm, and a policy that measures more
carefully still chooses it. The pairing itself does win against the
baselines that were in the grid --- FCFS by $6.5\%$, static partitioning
and EDF by about $3\%$, and at overload FCFS by $17\%$ --- which is why
the baseline set mattered more than the policy set, as the next
paragraph but one records.

\paragraph{The baseline that was missing.} For most of this project's
life the ``no partitioning'' arm was whole-die time-slicing between
tenants. Strict priority --- the whole die to the deadline-carrying
tenant whenever it has work, no partitioning ever, the obvious production
policy --- had never been run. It beats the adaptive scheduler on urgent
miss rate by $+134\%$ under arrivals and $+84.5\%$ under backlog,
intervals excluding zero under a seed cluster bootstrap, better in 15 of
20 cells and tied in 3; under arrivals it dominates, with identical video
goodput and identical Jain fairness. We then tested the two defects that
could have explained it --- a round barrier and a two-action space --- in
a 240-cell factorial, and neither does: every cell loses to priority by
$+83\%$ to $+434\%$, and six actions are far worse than two.

\textbf{How much of the grid could separate anything.} ``45 paired
configurations'' overstates the evidence, so we say what it is. The
probe and step-matched pairing produce \emph{identical} miss rates in 42
of the 45; three differ, all in the same direction. A bootstrap resample
that draws none of the three has probability $(42/45)^{45}=4.5\%$, above
$2.5\%$, so the interval's lower bound is zero as arithmetic and cannot
be read as significance --- the grid shows the probe does not help, not
that it hurts. The comparisons above are wider: pairing differs from
FCFS in 21 configurations, from static partitioning in 16, from EDF in
17. Two tenants are both runnable for only $26.1\%$ of the horizon at
load $0.6$ and $44.9\%$ at $1.05$, and three of the 45 configurations
draw no video tenant at all, so most of this grid could not have
separated any two partitioning policies whatever they did.

We state this as a negative because it was registered as one. The
comparison target was fixed by rule --- best baseline mean, computed not
chosen --- because the same method beats FCFS by $30\%$, the strongest
baseline by $22\%$, and its own ablation by $4.5\%$ on a single cell, and
which opponent a headline uses is the easiest thing to lose between a
table and a sentence.

\section{Three instruments that fitted artifacts}
\label{sec:method}

Reading a step's events one step late avoids draining the pipeline, but
the read only lands once the previous step's events have retired. When
the measured side's steps are short it never lands, and the variable
holds one stale value for an entire episode. The side whose steps are
short is the side whose CPU runs ahead, which in a mismatched pair is
exactly the side whose penalty is being measured.

The same defect produced three published cost models, each quantitative
and each wrong.

\textbf{One.} A two-episode ``serialisation transient'' in which each
side appeared to cost the sum of the two solo steps, fitting five splits
to within 3--9\% over a $3\times$ range of that sum. A real fit to a
stale variable.

\textbf{Two.} A wall clock introduced to replace the stale events, but
stopped without synchronising the device, so it measured the CPU's
submission rate: 14 steps of a 269\,ms quota read 156\,ms. The N-way
penalty of \S\ref{sec:peers} was withdrawn on that reading and restored
the next day when the clock was drained before it was stopped.

\textbf{Three.} The claim this paper's \S\ref{sec:mismatch} replaces:
mismatched tenants costing $1.00$--$1.06$ per side and Pareto-dominating
rotation at every split. Its own published text contains the refutation,
read at the time as a warm-up transient --- a first co-run episode of
$988$\,ms per step, ``exactly its own $183$\,ms solo step plus one whole
$805$\,ms CogVideoX step, with no spread at all''. No spread at all is
the signature of a variable that is not being updated. What we now
measure at that split is $5.987$, and equation~\ref{eq:pacing} is the
$988$\,ms account with a coefficient of $0.543$ instead of $1$ and as the
steady state rather than a transient.

\textbf{The control that catches all three} is cheap and was missing from
every one of them: after the co-run episodes, re-measure the solo with
the peers idle, then again after dropping the allocator cache, then again
with the peer streams destroyed. A penalty that does not vanish when the
peers stop is not a penalty. Applying it retroactively disqualifies the
comparator set we had intended to use for \S\ref{sec:mismatch} --- at
eight ways its post-episode solo reads $2048$\,ms against a pre-episode
$504$ --- which is why that section measures its own same-model baseline
in the same session.

What caught it was not reasoning. Running a second instrument --- events
bracketing the same step on the same stream, which physically
\emph{contain} the first --- over the same steps in the same process
gave $1.022$ where the first gave $10.888$. Containment made the
contradiction unarguable.

We record the working rules this produced, because the same class of
error --- a property of the measurement presenting itself as a property
of the thing measured --- occurred nine times: measure the null before
varying parameters; state the discriminating criterion and its direction
before running; a single co-run measurement is not evidence; the
measurement architecture must match the claimed deployment architecture;
and make the nuisance variable identical across arms rather than
trusting it to balance --- seed with position, deadline with arm, and
co-run state with arm each confounded a campaign before that was
routine.

\section{Threats to validity}

\textbf{Two SKUs, one vendor.} Results are from gfx1201 (RDNA4, 32
maskable units) and gfx90a (CDNA2, 104 per GCD). The two disagree in ways
that matter: bistability is present on the first and absent on the
second, and the aggregate solo throughput of $N$ slices peaks at four
ways on gfx90a where it is monotone on gfx1201, so the Amdahl form our
cost model assumes is refuted there outright. Nothing rules out that the
whole family of results depends on AMD's queueing and arbitration. No
number of additional cells on the same two cards substitutes for
cross-vendor agreement.

\textbf{Two models, one workpoint each.} SDXL at $768\times768$ and
CogVideoX-2b at $480\times720\times9$. Equation~\ref{eq:pacing}'s
coefficient is fitted across five splits of one pair; whether $k$ is a
constant of the hardware, of the step-length ratio, or of these two
models is not established. The obvious next measurement is a third model
with an intermediate step length.

\textbf{One process.} \S\ref{sec:peers} shows the mechanism does not
survive the process boundary, so every quantitative result here describes
a runtime that holds all tenants in one address space. For a deployment
of one server per tenant, the finding is the $24\%$ loss, not the table.

\textbf{Same-model co-location is not the primary workload.} The probe
results of \S\ref{sec:probe} are measured there because the mechanism is
inert on the mismatched one.

\textbf{Small $n$ on the slow state.} Ten processes. The interval
excludes zero; it is not tight.

\textbf{An unexplained ablation.} Making the runtime blind to the
externality term \emph{reduced} serial fallbacks where the mechanism and
a unit test both say it should increase them. Under investigation.

\section{Conclusion}

The claim this work set out to make --- that spatial partitioning raises
utilisation by a fixed factor --- is false as stated, and so was our
first replacement for it. The factor is not fixed, and it is not a
function of the split. It is drawn per process for same-model tenants,
it grows with the number of peers rather than with how busy the die is,
it vanishes across the process boundary, and for a mismatched pair it is
set by the slower tenant's step, so the faster tenant gets worse as you
give it more of the die.

That last one is why our scheduler lost. A partitioner needs
$\mathit{externality}(q, 32-q)$; this hardware supplies something whose
arguments include which models, how many, and in whose address space. Our
policies asked the first question, and every negative in this paper ---
the pre-registered grid, the probe, the loss to strict priority --- is
what that costs. A scheduler that wants to use spatial partitioning here
needs pair-specific measurements, an explicit refusal for configurations
it has none for, and time-slicing kept as a first-class action rather
than as the baseline it is assumed to beat.

We report our own retractions in full because three of them came from one
instrument defect and one control would have caught all three, and
because the version of this paper that would have been submitted a month
ago had a correct-looking quantitative account of a phenomenon that was
not happening.

\end{document}
